\documentclass[11pt]{article}

\usepackage[margin=1in]{geometry}
\usepackage{amsmath, amssymb, mathtools}
\usepackage{enumitem}
\usepackage{parskip}
\usepackage{xcolor}
\usepackage[colorlinks=true,linkcolor=blue,urlcolor=blue]{hyperref}

\title{CEC 315 --- Exam 3 Flashcards \\ \large Lectures 16--23}
\author{Signals and Systems (Spring 2026)}
\date{}

% Q/A environment
\newcommand{\Q}[1]{\par\medskip\noindent\textbf{Q:} #1}
\newcommand{\A}[1]{\par\noindent\textbf{A:} #1}
\newcommand{\Pit}[1]{\par\noindent\textit{Pitfall:} #1}

\begin{document}
\maketitle

\noindent\textbf{Coverage:} Laplace, Inverse Laplace, Unilateral Laplace, Z-Transform, Inverse Z, Unilateral Z, Sampling, Feedback Systems. \\
\textbf{Format:} Front/back Q\&A for rapid memorization drilling.

\tableofcontents

%============================================================
\section{Lecture 16 --- Laplace Transform and ROC}
%============================================================

\Q{Write the definition of the bilateral Laplace transform.}
\A{$X(s) = \int_{-\infty}^{\infty} x(t)\,e^{-st}\,dt$, where $s = \sigma + j\omega$. It is the Fourier transform of $x(t)e^{-\sigma t}$; the damping factor $e^{-\sigma t}$ tames signals whose FT diverges.}
\Pit{$X(s)$ alone is incomplete --- always state the ROC.}

\Q{Why do we need the Laplace transform when we already have the Fourier transform?}
\A{The FT fails for signals that are not absolutely integrable (e.g., $e^{2t}u(t)$). With $\sigma$ large enough, $e^{-\sigma t}$ forces convergence, so Laplace handles a strictly larger class of signals.}
\Pit{Laplace is \textbf{not} just ``FT with $\omega\to s$''; the substitution $s=j\omega$ works only if the $j\omega$-axis lies in the ROC.}

\Q{What is the Region of Convergence (ROC)?}
\A{The set of complex $s$ for which $\int x(t)e^{-st}dt$ converges. For rational $X(s)$, it is a vertical strip in the $s$-plane bounded by poles, parallel to the $j\omega$-axis.}
\Pit{ROC depends only on $\sigma=\operatorname{Re}\{s\}$, never on $\omega$.}

\Q{State the Laplace pair for $e^{-at}u(t)$.}
\A{$e^{-at}u(t) \overset{\mathcal{L}}{\longleftrightarrow} \dfrac{1}{s+a}$, ROC $\operatorname{Re}\{s\} > -\operatorname{Re}\{a\}$.}
\Pit{Same algebraic expression $1/(s+a)$ is also the transform of $-e^{-at}u(-t)$ with ROC $\operatorname{Re}\{s\} < -\operatorname{Re}\{a\}$ --- only the ROC distinguishes them.}

\Q{Laplace transforms of $\delta(t)$ and $u(t)$?}
\A{$\delta(t) \leftrightarrow 1$, ROC = all $s$. $u(t) \leftrightarrow 1/s$, ROC $\operatorname{Re}\{s\}>0$.}
\Pit{The pole of $u(t)$ sits on the $j\omega$-axis, so you cannot get the FT of $u(t)$ just by setting $s=j\omega$.}

\Q{Laplace transforms of $\cos(\omega_0 t)u(t)$ and $\sin(\omega_0 t)u(t)$?}
\A{$\cos\to\dfrac{s}{s^2+\omega_0^2}$, $\sin\to\dfrac{\omega_0}{s^2+\omega_0^2}$, both ROC $\operatorname{Re}\{s\}>0$.}
\Pit{Poles at $\pm j\omega_0$ are on the axis --- marginally stable, so FVT doesn't apply.}

\Q{Laplace pair for $\dfrac{t^{n-1}}{(n-1)!}e^{-at}u(t)$?}
\A{$\dfrac{t^{n-1}}{(n-1)!}e^{-at}u(t) \leftrightarrow \dfrac{1}{(s+a)^n}$, same ROC as the simple pole.}
\Pit{Repeated-pole numerators need the $(n-1)!$ normalization factor.}

\Q{ROC rule for a right-sided signal?}
\A{Right half-plane to the right of the rightmost pole: $\operatorname{Re}\{s\} > \max_i \operatorname{Re}\{p_i\}$.}
\Pit{``Right-sided'' means $x(t)=0$ for $t<T_0$; strictly causal ($T_0=0$) is the usual case.}

\Q{ROC rule for a left-sided signal?}
\A{Left half-plane to the left of the leftmost pole: $\operatorname{Re}\{s\} < \min_i \operatorname{Re}\{p_i\}$.}
\Pit{The ROC points \emph{away} from where the signal has its growing (``problematic'') behavior.}

\Q{ROC of a finite-duration signal?}
\A{The entire $s$-plane (possibly excluding $s=\pm\infty$) because the finite-limit integral always converges.}
\Pit{$\delta(t)$ and truncated signals have no poles.}

\Q{Can the ROC contain a pole?}
\A{No --- at a pole $X(s)\to\infty$, so the integral cannot converge there. Poles form the \emph{boundaries} of the ROC.}
\Pit{Never an interior point.}

\Q{When does $X(j\omega) = X(s)\big|_{s=j\omega}$?}
\A{If and only if the $j\omega$-axis lies inside the ROC, i.e., $\sigma = 0$ is a convergent line.}
\Pit{For $e^{2t}u(t)$, ROC is $\operatorname{Re}\{s\}>2$ --- no FT exists.}

\Q{Shape of the ROC for a two-sided signal?}
\A{A vertical strip between two consecutive poles: $\sigma_- < \operatorname{Re}\{s\} < \sigma_+$. Empty if the required bounds are inconsistent.}
\Pit{Always intersect the individual ROCs when using linearity.}

\Q{What symbols mark poles and zeros in a pole-zero plot?}
\A{$\times$ for poles (roots of denominator $D(s)$), $\circ$ for zeros (roots of numerator $N(s)$). $X(s) = K\dfrac{\prod(s-z_i)}{\prod(s-p_j)}$.}
\Pit{Zeros don't affect stability; only poles do.}

\Q{What is the Laplace transform of $\delta(t-t_0)$?}
\A{$\delta(t-t_0) \leftrightarrow e^{-st_0}$, ROC = all $s$. Finite duration $\Rightarrow$ no poles.}
\Pit{The shift shows up in the exponent, not in the ROC.}

\Q{How does Laplace relate to eigenfunctions of LTI systems?}
\A{$e^{st}$ is an eigenfunction: $y(t) = H(s)e^{st}$ where $H(s)=\int h(t)e^{-st}dt$ is the Laplace transform of the impulse response.}
\Pit{Valid only for LTI systems and for $s$ in the ROC of $H(s)$.}

%============================================================
\section{Lecture 17 --- Inverse Laplace and Properties}
%============================================================

\Q{Practical method for computing an inverse Laplace transform?}
\A{Partial-fraction expansion (PFE) + table lookup. Rewrite $X(s)$ as a sum of standard pairs; each term's direction (right- vs left-sided) is fixed by the ROC.}
\Pit{Partial fractions alone are ambiguous --- you need the ROC to pick the correct inverse per term.}

\Q{How do you handle a repeated pole of order $n$ at $s = p$ in PFE?}
\A{Include $n$ terms $\dfrac{B_1}{s-p} + \dfrac{B_2}{(s-p)^2} + \cdots + \dfrac{B_n}{(s-p)^n}$. Invert each with $\dfrac{t^{k-1}}{(k-1)!}e^{pt}u(t)$.}
\Pit{Forgetting any order is a common mistake --- list all orders up to $n$.}

\Q{How do you invert a pair of complex-conjugate poles?}
\A{Complete the square: $(s+a)^2 + \omega_d^2$. Use pairs $e^{-at}\cos(\omega_d t)u(t) \leftrightarrow \dfrac{s+a}{(s+a)^2+\omega_d^2}$ and $e^{-at}\sin(\omega_d t)u(t) \leftrightarrow \dfrac{\omega_d}{(s+a)^2+\omega_d^2}$.}
\Pit{Don't split into complex PFE terms --- use the damped sinusoid pairs.}

\Q{State the Laplace linearity property.}
\A{$ax(t)+by(t) \leftrightarrow aX(s)+bY(s)$. New ROC contains $\text{ROC}_x \cap \text{ROC}_y$.}
\Pit{If the two ROCs don't overlap, the sum has no Laplace transform.}

\Q{State the time-shift property.}
\A{$x(t-t_0) \leftrightarrow e^{-st_0}X(s)$, same ROC.}
\Pit{Algebraic $X(s)$ picks up $e^{-st_0}$; ROC unchanged. Not to be confused with $s$-shift.}

\Q{State the $s$-domain shift property.}
\A{$e^{s_0 t}x(t) \leftrightarrow X(s-s_0)$. ROC shifts by $\operatorname{Re}\{s_0\}$.}
\Pit{Time-shift multiplies $X(s)$; $s$-shift changes its argument --- they are duals, not the same.}

\Q{State the time-scaling property.}
\A{$x(at) \leftrightarrow \dfrac{1}{|a|}X(s/a)$. ROC scales by $|a|$.}
\Pit{$a<0$ flips the ROC across the $j\omega$-axis.}

\Q{State the differentiation-in-time property (bilateral).}
\A{$\dfrac{dx}{dt} \leftrightarrow sX(s)$; ROC contains the original ROC.}
\Pit{No initial-condition term --- that's unilateral only.}

\Q{State the integration-in-time property.}
\A{$\displaystyle\int_{-\infty}^{t} x(\tau)d\tau \leftrightarrow \dfrac{X(s)}{s}$; ROC $= \text{ROC}_x \cap \{\operatorname{Re}\{s\}>0\}$.}
\Pit{Dividing by $s$ adds a pole at the origin and may shrink the ROC.}

\Q{State the differentiation-in-$s$ property.}
\A{$-tx(t) \leftrightarrow \dfrac{dX(s)}{ds}$, equivalently $tx(t) \leftrightarrow -X'(s)$.}
\Pit{Watch the minus sign.}

\Q{State the convolution property.}
\A{$x(t)*y(t) \leftrightarrow X(s)Y(s)$; ROC contains $\text{ROC}_x\cap\text{ROC}_y$.}
\Pit{This is why Laplace turns differential equations into algebra.}

\Q{State the Initial Value Theorem (IVT).}
\A{If $x(t)=0$ for $t<0$ and $x$ has no impulse at the origin: $x(0^+) = \lim_{s\to\infty} sX(s)$.}
\Pit{Requires $x(t)$ causal and regular at $t=0$; impulses break it.}

\Q{State the Final Value Theorem (FVT).}
\A{$x(\infty) = \lim_{s\to 0} sX(s)$, \textbf{valid only if} all poles of $sX(s)$ are in the open LHP (one allowed at $s=0$).}
\Pit{Applying FVT to oscillating or unstable signals gives nonsense.}

\Q{How do you compute the residue at a simple pole $p$?}
\A{$B = \lim_{s\to p}(s-p)X(s)$. That term becomes $\dfrac{B}{s-p}$.}
\Pit{Valid only for simple poles; repeated poles require derivatives.}

\Q{Residues at a repeated pole $(s-p)^n$?}
\A{$B_k = \dfrac{1}{(n-k)!}\left.\dfrac{d^{n-k}}{ds^{n-k}}\left[(s-p)^n X(s)\right]\right|_{s=p}$ for $k=1,\dots,n$.}
\Pit{Factorial and differentiation order --- off-by-one errors are common.}

\Q{How do you evaluate $|H(j\omega)|$ geometrically from a pole-zero plot?}
\A{$|H(j\omega_0)| = K\dfrac{\prod \text{(distances to zeros)}}{\prod \text{(distances to poles)}}$. Phase $=$ sum of zero angles $-$ sum of pole angles.}
\Pit{A pole near the $j\omega$-axis produces a peak; a zero near the axis produces a notch.}

\Q{Laplace transform of $e^{-at}\cos(\omega_d t)u(t)$?}
\A{$\dfrac{s+a}{(s+a)^2+\omega_d^2}$, ROC $\operatorname{Re}\{s\}>-a$.}
\Pit{Numerator is $s+a$, not $s$.}

%============================================================
\section{Lecture 18 --- System Analysis via Unilateral Laplace}
%============================================================

\Q{Given an LCCDE $\sum a_k y^{(k)}(t) = \sum b_k x^{(k)}(t)$, how do you find $H(s)$?}
\A{Take Laplace transforms with \textbf{zero} ICs, substitute $d^k/dt^k \to s^k$: $H(s) = \dfrac{Y(s)}{X(s)} = \dfrac{\sum b_k s^k}{\sum a_k s^k}$.}
\Pit{ICs must be zero for the transfer function; nonzero ICs give ZIR contributions, not $H(s)$.}

\Q{What does the ROC of $H(s)$ tell you about causality?}
\A{The system is causal iff $H(s)$ is rational, proper ($M\le N$), and the ROC is a right half-plane to the right of the rightmost pole.}
\Pit{Causality is an ROC statement, not just a pole statement.}

\Q{What does the ROC tell you about BIBO stability?}
\A{BIBO stable iff the ROC contains the $j\omega$-axis (so $h(t)$ is absolutely integrable).}
\Pit{Stability is about the impulse response, not poles alone.}

\Q{State the ``Golden Rule'' for causal LTI systems.}
\A{Causal LTI $\Rightarrow$ BIBO stable iff \emph{all} poles of $H(s)$ are strictly in the open LHP ($\operatorname{Re}\{p_i\}<0$).}
\Pit{Poles on the $j\omega$-axis $\Rightarrow$ marginally stable, \textbf{not} BIBO stable.}

\Q{Block-diagram reduction rules?}
\A{Cascade: $H_1 H_2$. Parallel: $H_1 + H_2$. Negative feedback (forward $G$, feedback $F$): $\dfrac{G}{1+GF}$.}
\Pit{Positive feedback uses $1-GF$; sign matters.}

\Q{Define the unilateral Laplace transform.}
\A{$\mathcal{X}(s) = \displaystyle\int_{0^-}^{\infty} x(t)e^{-st}\,dt$. The $0^-$ lower limit captures impulses and ICs at $t=0$.}
\Pit{Always use $0^-$ vs $0^+$ if $x(t)$ has a jump or impulse at $t=0$.}

\Q{Unilateral differentiation rule (first derivative)?}
\A{$\mathcal{L}_u\{x'(t)\} = sX(s) - x(0^-)$.}
\Pit{Sign is \textbf{minus}; contrast with the $+$ in the unilateral Z shift rule.}

\Q{Unilateral differentiation rule (second derivative)?}
\A{$\mathcal{L}_u\{x''(t)\} = s^2 X(s) - sx(0^-) - x'(0^-)$.}
\Pit{Order matters: $sx(0^-)$ (not $x(0^-)s^2$), $x'(0^-)$ is the final term.}

\Q{How do you solve an LCCDE with nonzero ICs using unilateral Laplace?}
\A{Transform both sides, apply IC-bearing derivative formulas, solve for $Y(s)$, and invert. Output $=$ ZSR $+$ ZIR.}
\Pit{Don't forget IC terms for each derivative.}

\Q{What is the zero-state response (ZSR)?}
\A{Response with all ICs zero: $Y_{\text{ZSR}}(s) = H(s)X(s)$.}
\Pit{Not the full answer if ICs are nonzero.}

\Q{What is the zero-input response (ZIR)?}
\A{Response to ICs alone (input set to zero); composed entirely of natural modes of the homogeneous equation.}
\Pit{ZIR lives on poles of $H(s)$, independent of $X(s)$.}

\Q{How do you distinguish natural- vs forced-response modes?}
\A{Natural modes $\leftarrow$ poles of $H(s)$. Forced modes $\leftarrow$ poles of $X(s)$.}
\Pit{If they coincide, you get $t e^{pt}$-style terms.}

\Q{For $y'' + 3y' + 2y = x(t)$, what is $H(s)$?}
\A{$H(s) = \dfrac{1}{s^2+3s+2} = \dfrac{1}{(s+1)(s+2)}$. Both poles in LHP $\Rightarrow$ causal $+$ stable.}
\Pit{Factor to read off pole locations.}

\Q{For a causal LTI system with $H(s) = \dfrac{1}{s-1}$, is it stable?}
\A{No --- pole at $s=1$ is in the RHP. Golden Rule fails.}
\Pit{Causality alone doesn't imply stability.}

\Q{Feedback transfer function with forward $G(s)$ and feedback $F(s)$ (negative feedback)?}
\A{$\dfrac{Y(s)}{X(s)} = \dfrac{G(s)}{1+G(s)F(s)}$.}
\Pit{Unity feedback $F=1$: $Y/X = G/(1+G)$. Don't drop the $+1$.}

\Q{Why does causal $+$ stable imply the rightmost pole has $\operatorname{Re}<0$?}
\A{Causal $\Rightarrow$ ROC is right of rightmost pole. Stable $\Rightarrow j\omega$-axis in ROC. Both $\Rightarrow$ rightmost pole in open LHP.}
\Pit{Combines into ``all poles in open LHP.''}

%============================================================
\section{Lecture 19 --- Z-Transform and ROC}
%============================================================

\Q{Write the definition of the bilateral Z-transform.}
\A{$X(z) = \displaystyle\sum_{n=-\infty}^{\infty} x[n]z^{-n}$, with $z = re^{j\omega}$.}
\Pit{It is the DTFT of $x[n]r^{-n}$; when $r=1$, Z reduces to DTFT.}

\Q{What replaces the $j\omega$-axis in the $z$-plane?}
\A{The unit circle $|z|=1$. DTFT: $X(e^{j\omega}) = X(z)\big|_{z=e^{j\omega}}$.}
\Pit{Valid only if the unit circle $\subset$ ROC.}

\Q{Z-transform pair for $a^n u[n]$?}
\A{$a^n u[n] \leftrightarrow \dfrac{1}{1-az^{-1}} = \dfrac{z}{z-a}$, ROC $|z|>|a|$. Pole at $z=a$.}
\Pit{The pole is at $z=a$, not $z=1/a$.}

\Q{Z-transform pair for $-a^n u[-n-1]$?}
\A{$-a^n u[-n-1] \leftrightarrow \dfrac{1}{1-az^{-1}}$, ROC $|z|<|a|$.}
\Pit{Same algebraic $X(z)$ as $a^n u[n]$ --- ROC distinguishes them.}

\Q{Z-transform of $\delta[n]$ and $\delta[n-k]$?}
\A{$\delta[n]\leftrightarrow 1$, all $z$. $\delta[n-k]\leftrightarrow z^{-k}$, $|z|>0$ ($k>0$) or $|z|<\infty$ ($k<0$).}
\Pit{Shift introduces an exclusion at $0$ or $\infty$.}

\Q{Z-transform of $u[n]$?}
\A{$u[n]\leftrightarrow \dfrac{1}{1-z^{-1}} = \dfrac{z}{z-1}$, ROC $|z|>1$.}
\Pit{Pole at $z=1$ is on the unit circle, so DTFT doesn't follow from plain substitution.}

\Q{Shape of a Z-transform ROC?}
\A{An annular ring $r_1 < |z| < r_2$ centered at the origin, never containing a pole.}
\Pit{Depends only on $|z|$, not $\arg z$.}

\Q{ROC rule for a right-sided (causal) sequence?}
\A{$|z| > |p_{\max}|$, exterior of the circle through the outermost pole. May extend to $\infty$.}
\Pit{Truly causal sequences must include $z=\infty$ in the ROC.}

\Q{ROC rule for a left-sided sequence?}
\A{$|z|<|p_{\min}|$, interior of the circle through the innermost pole.}
\Pit{``Inside'' means smaller $|z|$, i.e., closer to the origin.}

\Q{ROC for a two-sided sequence?}
\A{Annular ring $|p_k|<|z|<|p_{k+1}|$ between two consecutive poles ordered by magnitude. May be empty.}
\Pit{Individual ROCs must intersect.}

\Q{ROC of a finite-length sequence?}
\A{Entire $z$-plane, excluding possibly $z=0$ (positive-index samples) and/or $z=\infty$ (negative-index samples).}
\Pit{Causal finite sequences exclude only $z=0$.}

\Q{When does the DTFT exist (via ROC)?}
\A{Iff the unit circle $|z|=1$ lies in the ROC.}
\Pit{Pole on the unit circle $\Rightarrow$ no ordinary DTFT.}

\Q{CT--DT parallel: what maps to what?}
\A{$s$-plane LHP $\leftrightarrow$ unit-circle interior. $j\omega$-axis $\leftrightarrow$ unit circle. RHP $\leftrightarrow$ exterior.}
\Pit{The literal map is $z=e^{sT}$; the parallel is conceptual.}

\Q{How does pole magnitude relate to sequence growth?}
\A{Causal: $|p|<1 \Rightarrow$ decay; $|p|=1 \Rightarrow$ constant amplitude; $|p|>1 \Rightarrow$ growth.}
\Pit{Sign of a real pole controls alternation (e.g., $(-0.5)^n$ oscillates while decaying).}

\Q{Z-transform of $na^n u[n]$?}
\A{$na^n u[n]\leftrightarrow \dfrac{az^{-1}}{(1-az^{-1})^2}$, ROC $|z|>|a|$.}
\Pit{Derived from $nx[n]\leftrightarrow -z\,dX/dz$.}

\Q{Why is Z the DT analog of Laplace?}
\A{Both generalize FT/DTFT by introducing a damping/growth factor (Laplace: $e^{-\sigma t}$; Z: $r^{-n}$). Same structure in pairs, ROC, stability.}
\Pit{Stability region differs: LHP vs unit-circle interior.}

%============================================================
\section{Lecture 20 --- Inverse Z-Transform and Properties}
%============================================================

\Q{Standard method for inverse Z?}
\A{Partial-fraction expansion \textbf{in $z^{-1}$} and table lookup; ROC determines the sign of each term.}
\Pit{Work in $z^{-1}$, not $z$, or the standard form $1/(1-az^{-1})$ won't appear.}

\Q{ROC \emph{outside} a pole $\Rightarrow$ which sign?}
\A{Right-sided: $\dfrac{A}{1-az^{-1}} \rightarrow A\,a^n u[n]$.}
\Pit{Only valid if that pole is on the inside of the ROC ring.}

\Q{ROC \emph{inside} a pole $\Rightarrow$ which sign?}
\A{Left-sided: $\dfrac{A}{1-az^{-1}} \rightarrow -A\,a^n u[-n-1]$.}
\Pit{Forgetting the minus sign and the $-1$ shift in $u[-n-1]$ is common.}

\Q{Inverse for complex-conjugate poles at $re^{\pm j\omega_0}$?}
\A{$r^n\cos(\omega_0 n)u[n]$ or $r^n\sin(\omega_0 n)u[n]$, ROC $|z|>r$.}
\Pit{Use real pairs; don't do complex PFE.}

\Q{Z linearity?}
\A{$ax_1[n]+bx_2[n]\leftrightarrow aX_1(z)+bX_2(z)$, ROC contains $\text{ROC}_1\cap\text{ROC}_2$.}
\Pit{Non-overlapping ROCs $\Rightarrow$ no Z-transform of the sum.}

\Q{Z time shift?}
\A{$x[n-n_0]\leftrightarrow z^{-n_0}X(z)$, ROC unchanged (except maybe at $z=0$ or $\infty$).}
\Pit{$z^{-1}$ is the unit delay operator.}

\Q{Z-domain scaling?}
\A{$z_0^n x[n]\leftrightarrow X(z/z_0)$. ROC bounds multiply by $|z_0|$.}
\Pit{Special case $(-1)^n x[n]\leftrightarrow X(-z)$.}

\Q{Z convolution?}
\A{$x_1[n]*x_2[n]\leftrightarrow X_1(z)X_2(z)$.}
\Pit{Why $H(z)=Y(z)/X(z)$ is meaningful for LTI systems.}

\Q{$z$-differentiation?}
\A{$nx[n]\leftrightarrow -z\,\dfrac{dX(z)}{dz}$.}
\Pit{The $-z$ multiplier, not just $-d/dz$.}

\Q{Time reversal?}
\A{$x[-n]\leftrightarrow X(z^{-1})$; ROC inverts: $r_1<|z|<r_2 \to 1/r_2<|z|<1/r_1$.}
\Pit{Causal reversed becomes anti-causal; ROC flips accordingly.}

\Q{First difference and accumulation properties?}
\A{$x[n]-x[n-1]\leftrightarrow (1-z^{-1})X(z)$. $\displaystyle\sum_{k\le n} x[k]\leftrightarrow \dfrac{X(z)}{1-z^{-1}}$.}
\Pit{Accumulation adds a pole at $z=1$.}

\Q{Z IVT (causal)?}
\A{$x[0] = \lim_{z\to\infty} X(z)$, valid for causal $x[n]$.}
\Pit{Causality required.}

\Q{Z FVT?}
\A{$x[\infty] = \lim_{z\to 1}(1-z^{-1})X(z)$, valid only if all poles of $(1-z^{-1})X(z)$ are strictly inside the unit circle (one allowed at $z=1$).}
\Pit{Check pole locations before applying.}

\Q{Long-division method for inverse Z?}
\A{Polynomial long division of $X(z)$ in $z^{-1}$ yields $x[0]+x[1]z^{-1}+x[2]z^{-2}+\cdots$; coefficients are sample values.}
\Pit{Practical only for the first few samples; direction depends on whether you divide in $z^{-1}$ or $z$.}

\Q{Residue at a simple pole $p$ (in $z^{-1}$ form)?}
\A{$A_k = (1-p_k z^{-1})X(z)\big|_{z=p_k}$.}
\Pit{Using $z$-form yields residues with an extra factor of $z$.}

\Q{Meaning of a $z^{-1}$ block in a DSP block diagram?}
\A{Unit delay (one-sample memory element). Multiplying by $z^{-1}$ in the Z-domain $\equiv$ delay by one sample.}
\Pit{Fundamental building block of FIR/IIR filters.}

%============================================================
\section{Lecture 21 --- System Analysis via Unilateral Z}
%============================================================

\Q{How do you derive $H(z)$ from a difference equation?}
\A{Zero ICs, $y[n-k]\to z^{-k}Y(z)$, $x[n-k]\to z^{-k}X(z)$, solve for $Y/X$.}
\Pit{Must use zero ICs for the transfer function.}

\Q{For $y[n]-0.5y[n-1]=x[n]$, what is $H(z)$?}
\A{$H(z) = \dfrac{1}{1-0.5z^{-1}}$. Pole at $z=0.5$ inside unit circle $\Rightarrow$ causal $+$ stable.}
\Pit{Magnitude of the pole matters, not the sign.}

\Q{Golden Rule for DT causal LTI stability.}
\A{BIBO stable iff all poles of $H(z)$ satisfy $|p_i|<1$.}
\Pit{Poles on the unit circle $\Rightarrow$ marginal stability only.}

\Q{Causality implies what about the ROC of $H(z)$?}
\A{Exterior of a circle: $|z|>|p_{\max}|$, extending to $\infty$; rational with $M\le N$.}
\Pit{Causality (ROC) and stability (unit circle $\subset$ ROC) are separate.}

\Q{Define the unilateral Z-transform.}
\A{$\mathcal{X}(z) = \displaystyle\sum_{n=0}^{\infty} x[n]z^{-n}$.}
\Pit{Ignores $n<0$; ICs enter only via shift rules.}

\Q{Unilateral Z shift rule for $x[n-1]$?}
\A{$x[n-1]\leftrightarrow z^{-1}X(z) + x[-1]$. Note the \textbf{plus}.}
\Pit{Opposite sign from unilateral Laplace differentiation.}

\Q{Unilateral Z shift rule for $x[n-2]$?}
\A{$x[n-2]\leftrightarrow z^{-2}X(z) + x[-1]z^{-1} + x[-2]$.}
\Pit{$x[-1]$ multiplies $z^{-1}$; $x[-2]$ stands alone.}

\Q{Step response in Z-domain?}
\A{$Y_{\text{step}}(z) = H(z)U(z) = H(z)\cdot\dfrac{1}{1-z^{-1}}$, then invert.}
\Pit{Step response has an added pole at $z=1$; FVT applies only if system is otherwise stable.}

\Q{ZSR vs ZIR decomposition (DT)?}
\A{$Y = Y_{\text{ZSR}} + Y_{\text{ZIR}}$. ZSR from $H(z)X(z)$; ZIR from ICs with $X(z)=0$.}
\Pit{Same decomposition as CT unilateral Laplace.}

\Q{What are natural modes of a DT system?}
\A{Terms $p_k^n$ from poles of $H(z)$; repeated poles give $n p_k^n$, etc.}
\Pit{These are the system's free-response eigenmodes.}

\Q{DT analog of an $RC$ time constant?}
\A{For a real pole $0<p<1$, time constant $= -1/\ln|p|$ samples.}
\Pit{Only meaningful for $|p|<1$.}

\Q{Can an FIR filter be unstable?}
\A{No --- FIR filters have poles only at $z=0$ (or $\infty$), always BIBO stable.}
\Pit{FIR design effort is in shaping the zeros.}

\Q{Pole-zero condition for real-valued impulse response?}
\A{Poles and zeros in complex-conjugate pairs.}
\Pit{Same rule as CT.}

\Q{Steady-state sinusoidal response of a stable DT system?}
\A{Evaluate $H(e^{j\omega_0})$; magnitude and phase give output amplitude and phase shift.}
\Pit{Only if the unit circle $\subset$ ROC.}

\Q{ZIR for $y[n]+0.9y[n-1]=x[n]$ with $y[-1]=2$, $x[n]=0$?}
\A{$Y(z)+0.9(z^{-1}Y(z)+y[-1])=0\Rightarrow Y(z)(1+0.9z^{-1}) = -1.8\Rightarrow y[n] = -1.8(-0.9)^n u[n]$.}
\Pit{Apply the unilateral shift rule carefully; sign of the IC on each side matters.}

\Q{Open-loop vs closed-loop poles in a DT feedback system?}
\A{Open-loop: poles of $G(z)H(z)$. Closed-loop: roots of $1+G(z)H(z)=0$.}
\Pit{Feedback can move poles in or out of the unit circle.}

%============================================================
\section{Lecture 22 --- Sampling}
%============================================================

\Q{State the Shannon sampling theorem.}
\A{A band-limited signal with $X(j\omega)=0$ for $|\omega|>\omega_M$ can be perfectly reconstructed from samples spaced $T$ apart provided $\omega_s = 2\pi/T > 2\omega_M$ (strict).}
\Pit{$\omega_s = 2\omega_M$ is \textbf{not} sufficient in general.}

\Q{Define the Nyquist rate.}
\A{$\omega_{\text{Nyquist}} = 2\omega_M$ (rad/s) or $f_{\text{Nyquist}} = 2f_M$ (Hz).}
\Pit{Nyquist \emph{rate} is a property of the signal; Nyquist \emph{frequency} ($\omega_s/2$) is a property of the sampler.}

\Q{Impulse train $p(t)$ and its role in sampling?}
\A{$p(t) = \sum_n \delta(t-nT)$, so $x_p(t) = x(t)p(t) = \sum_n x(nT)\delta(t-nT)$.}
\Pit{Multiplication in time $\Rightarrow$ convolution in frequency $\Rightarrow$ spectral replicas.}

\Q{Frequency-domain relation for impulse-train sampling?}
\A{$X_p(j\omega) = \dfrac{1}{T}\displaystyle\sum_{k=-\infty}^{\infty} X(j(\omega-k\omega_s))$.}
\Pit{The $1/T$ scaling must be cancelled by a gain $T$ in the reconstruction filter.}

\Q{Describe ideal reconstruction.}
\A{Pass $x_p(t)$ through an ideal LPF with gain $T$ and cutoff $\omega_c$ with $\omega_M < \omega_c < \omega_s-\omega_M$.}
\Pit{Gain $T$, not $1$.}

\Q{Write the band-limited (sinc) interpolation formula.}
\A{$x(t) = \displaystyle\sum_{n=-\infty}^{\infty} x(nT)\,\text{sinc}\!\left(\dfrac{t-nT}{T}\right)$, $\text{sinc}(u)=\sin(\pi u)/(\pi u)$.}
\Pit{Sinc kernels have infinite support --- practical systems truncate.}

\Q{What is aliasing and when does it occur?}
\A{If $\omega_s<2\omega_M$, spectral replicas overlap and fold back into the baseband. Aliased frequency: $f_a = |f_s-f_0|$.}
\Pit{Aliasing is \textbf{irreversible} --- anti-alias filter must be applied before sampling.}

\Q{Relation between DT and CT frequencies?}
\A{$\Omega = \omega T$. DT band $[-\pi,\pi]$ corresponds to CT $[-\omega_s/2,\omega_s/2]$.}
\Pit{Watch the units: rad/sample vs rad/s.}

\Q{What is a zero-order hold?}
\A{$h_{\text{ZOH}}(t)=1$ for $0\le t<T$, 0 else. $H(j\omega) = T\,\text{sinc}(\omega T/2\pi)e^{-j\omega T/2}$.}
\Pit{Introduces amplitude droop and a half-sample delay.}

\Q{What is a first-order hold?}
\A{Linear interpolation between samples; triangular impulse response of length $2T$.}
\Pit{Better than ZOH but still not ideal.}

\Q{Nyquist rates of $x(t-a)$, $x(2t)$, $x^2(t)$?}
\A{$x(t-a)$: unchanged. $x(2t)$: doubled. $x^2(t)$: doubled (convolution doubles spectral support).}
\Pit{Squaring/multiplying doubles bandwidth.}

\Q{What happens if you sample a 1 kHz cosine at 1.5 kHz?}
\A{$f_s/2=750$ Hz, so 1 kHz aliases to $|1000-1500|=500$ Hz.}
\Pit{Produces wagon-wheel / moir\'e effects.}

\Q{Why does the reconstruction LPF need gain $T$?}
\A{Because the sampling operation scales each spectral replica by $1/T$.}
\Pit{Unit check: $T\cdot(1/T)=1$.}

\Q{DSP of continuous-time signals --- block chain?}
\A{CT anti-aliasing LPF $\to$ C/D $\to$ DT system $H(e^{j\Omega}) \to$ D/C (reconstruction LPF). Effective CT transfer: $H_{\text{eff}}(j\omega) = H(e^{j\omega T})$ for $|\omega|<\pi/T$.}
\Pit{Valid only within the Nyquist band.}

\Q{Why is the anti-aliasing filter before, not after, sampling?}
\A{Because aliasing is irreversible --- folded content cannot be separated from true baseband content after sampling.}
\Pit{Filtering afterward only removes (already corrupted) baseband.}

\Q{Sampling theorem in Hz form?}
\A{$f_s > 2f_M$, i.e., $T < 1/(2f_M)$.}
\Pit{$\omega$ is rad/s; $f$ is Hz; they differ by $2\pi$.}

%============================================================
\section{Lecture 23 --- Linear Feedback Systems}
%============================================================

\Q{Closed-loop transfer function for unity feedback with forward gain $H(s)$?}
\A{$Q(s) = \dfrac{H(s)}{1+H(s)}$.}
\Pit{Unity case; non-unity uses $H/(1+HF)$.}

\Q{Closed-loop transfer function for non-unity negative feedback ($G$ forward, $H$ feedback)?}
\A{$Q(s) = \dfrac{G(s)}{1 + G(s)H(s)}$.}
\Pit{Positive feedback uses $1-GH$.}

\Q{Where are the closed-loop poles?}
\A{Roots of the characteristic equation $1 + G(s)H(s) = 0$, i.e., $GH = -1$.}
\Pit{Closed-loop poles $\ne$ open-loop poles of $GH$.}

\Q{What is the sensitivity function?}
\A{$S(s) = \dfrac{1}{1+L(s)}$ where $L = GH$. Large $|L| \Rightarrow$ low sensitivity to plant variation and disturbances.}
\Pit{Very high gain also amplifies noise and can destabilize.}

\Q{State the simplified Nyquist criterion (open-loop stable).}
\A{Closed-loop stable iff the Nyquist plot of $L(j\omega)$ does not encircle $-1/K$ (or $-1$ if $K$ absorbed into $L$).}
\Pit{General form counts encirclements; simplified version assumes open-loop stability.}

\Q{Define gain margin (GM).}
\A{At $\omega_1$ where $\angle L(j\omega_1) = -180^\circ$, $\text{GM} = -20\log_{10}|L(j\omega_1)|$ dB.}
\Pit{A first-order $L(s)$ has infinite GM (never reaches $-180^\circ$).}

\Q{Define phase margin (PM).}
\A{At gain crossover $\omega_2$ where $|L(j\omega_2)|=1$, $\text{PM} = 180^\circ + \angle L(j\omega_2)$.}
\Pit{Typical target: $45^\circ$--$60^\circ$.}

\Q{When is a closed-loop stable from Bode plots alone?}
\A{Both GM and PM positive (assuming open-loop stable, single crossover).}
\Pit{Multiple crossovers require Nyquist.}

\Q{Can feedback stabilize an unstable plant?}
\A{Yes --- feedback relocates poles; sufficient loop gain can push RHP poles into the LHP. But excessive gain can do the opposite.}
\Pit{Stability is nonmonotonic in gain.}

\Q{Maximum tolerable time delay given PM at $\omega_2$?}
\A{$\tau_{\max} = \dfrac{\text{PM (in radians)}}{\omega_2}$.}
\Pit{Convert PM from degrees to radians first.}

\Q{Cascade and parallel block rules?}
\A{Cascade: $H_1 H_2$. Parallel: $H_1 + H_2$.}
\Pit{Loading effects may invalidate cascade in physical systems.}

\Q{What is the loop gain and why is it important?}
\A{$L(s) = G(s)H(s)$. Determines sensitivity, stability margins, and closed-loop poles via $1+L=0$.}
\Pit{Define break points and signs consistently.}

\Q{How do you sketch the Nyquist plot?}
\A{Plot $L(j\omega)$ for $\omega:0\to\infty$, reflect for $\omega:-\infty\to 0$, close with a large semicircle, and count encirclements of $-1$.}
\Pit{Poles on the $j\omega$-axis require small indentations.}

\Q{Effect of increasing forward gain $K$ on the Nyquist plot?}
\A{Scales the plot radially by $K$; equivalently, count encirclements of $-1/K$.}
\Pit{Large $K$ can push the plot to enclose $-1$, destabilizing the loop.}

\Q{Open-loop vs closed-loop distinction?}
\A{Open loop: $Y=GX$. Closed loop: $Y = GX - GHY \Rightarrow Y = GX/(1+GH)$.}
\Pit{Open loop relies on accurate plant knowledge; closed loop tolerates uncertainty.}

\Q{Bode stability rule of thumb for stable minimum-phase $L(s)$?}
\A{If $|L|$ crosses $0$ dB with slope $\approx -20$ dB/dec, PM is usually adequate; $-40$ dB/dec often indicates marginal/unstable.}
\Pit{Heuristic only --- confirm with exact GM/PM or Nyquist.}

\Q{What role does feedback play in disturbance rejection?}
\A{Output-referred disturbances are attenuated by $1/(1+L)$; high $|L|$ at disturbance frequencies gives good rejection.}
\Pit{Rejection is frequency-dependent; shape $L(j\omega)$ accordingly.}

\bigskip
\begin{center}
\textit{Prepared for CEC 315 Exam 3 --- Spring 2026. Drill until the questions feel boring, then take the exam.}
\end{center}

\end{document}
